\documentclass{situationist-preprint}

\title{Why Unambiguous Revelation Would Constrain Freedom}
\msid{REL-03-b}
\msfamily{REL-03}
\msdiscipline{Philosophy of Religion}
\mstype{Consequence}
\mskeywords{divine hiddenness; revelation; freedom of response; signalling games; coercion}

\begin{document}
\maketitle

\begin{abstract}
REL-03-a argues that hiddenness is the equilibrium a perfectly loving God would choose, because unambiguous revelation would constrain the freedom of creatures' response. We derive that premise in general form from the structure of the hiddenness game. Whenever revelation raises the prudential stakes of turning toward God above every type's intrinsic preference for turning away, the best response no longer varies with the creature's disposition. The constraint is not coercion and does not remove free agency in the sense of reasons-responsiveness; what it removes is the dependence of the response on who the creature is, which is what makes the response worth having.
\end{abstract}

\section{The result}

In the hiddenness game of REL-03-a (Definition~2.1), God chooses the clarity $c$ of a signal, and creatures of two types respond by turning toward God or away. The response is \emph{free} at $c$ if the best responses of the two types differ (Definition~2.4). Lemma~3.1 of REL-03-a shows, for the payoffs of Definition~2.2, that under unambiguous revelation both types turn toward God, so that the response is not free. We state the underlying structure without the particular numbers.

Let each type $\tau$ receive $a_\tau + s(c)$ for turning toward God and $b_\tau$ for turning away, where $a_\tau$ and $b_\tau$ are the intrinsic values of the two responses to that type and $s(c)$ is the prudential component, increasing in clarity. Let $\tau$ be \emph{disposed toward} God if $a_\tau > b_\tau$ and \emph{disposed away} otherwise.

\begin{proposition}\label{p:pool}
If at clarity $c$ the prudential component exceeds $b_\tau - a_\tau$ for every type $\tau$, every type's best response at $c$ is to turn toward God, and the response at $c$ is not free whenever some type is disposed away.
\end{proposition}

\begin{proof}
For each $\tau$, $a_\tau + s(c) > b_\tau$, so turning toward God is the best response of every type, and the best responses coincide. A type disposed away would turn away at $s(c) = 0$, so the response at $c$ differs from what that type's disposition alone would yield.
\end{proof}

Assumption~2.3 of REL-03-a (Stakes) says that unambiguous revelation meets the hypothesis of Proposition~\ref{p:pool} and hiddenness does not. With the payoffs of REL-03-a, $a = 2$, $b = 0$ for the open type and $a = 0$, $b = 3$ for the closed type, with $s(c) = c$; at $c = 4$ the closed type receives $4 > 3$, and Lemma~3.1 of REL-03-a is the instance. The proposition shows that the result does not depend on the numbers: whatever the intrinsic values, a signal clear enough to make the prudential stakes dominate makes the response independent of disposition.

\section{What is constrained}

Proposition~\ref{p:pool} does not say that revealed creatures act unfreely in the sense of ACT-01. Each acts from its own reasons-responsive mechanism, and a creature disposed away who turns toward God under revelation does so for reasons, namely prudential ones. What is lost is the dependence of the response on the creature's orientation. Under hiddenness the response reveals who the creature is; under revelation it reveals only that the stakes are known. The value a perfectly loving God places on a response (Definition~2.5 and Assumption~2.6 of REL-03-a) is the value of a response that expresses the creature's orientation, as the gratitude or love of a person is valuable because it is given rather than extracted \citep{Hick1966}.

\begin{objection}
This is the coercion argument, and revelation need not threaten anyone \citep{Murray1993}.
\end{objection}

\begin{reply}
Proposition~\ref{p:pool} requires no threat. The prudential component may consist entirely of benefits: the goods of relationship with a being one knows to exist and to be perfectly good. It suffices that, once these are known with certainty, they outweigh every type's intrinsic preference. The constraint is structural, a pooling of types, and it holds however benign the stakes.
\end{reply}

\begin{objection}
Schellenberg's claim concerns being in a position to believe, not being motivated to respond \citep{Schellenberg1993}.
\end{objection}

\begin{reply}
In the game, belief and stakes come together: the clarity of the signal both positions the creature to believe and fixes the prudential component, since the prudential reasons are reasons only for one who believes. REL-03-a accordingly counts the creature's being in a position to believe as a good, and Theorem~3.3 of REL-03-a shows that God gives as much of it as is compatible with a free response.
\end{reply}

\begin{thebibliography}{99}\small
\bibitem[Hick(1966)]{Hick1966} Hick, J. (1966). \emph{Evil and the God of Love}. Macmillan.
\bibitem[Murray(1993)]{Murray1993} Murray, M.~J. (1993). Coercion and the hiddenness of God. \emph{American Philosophical Quarterly}, 30(1), 27--38.
\bibitem[Schellenberg(1993)]{Schellenberg1993} Schellenberg, J.~L. (1993). \emph{Divine Hiddenness and Human Reason}. Cornell University Press.
\end{thebibliography}

\end{document}
